%% 01-problem.tex — The gerrymandering problem

\section{The Problem}

\noindent\textbf{Every ten years, after the Census, the party in power in most
state legislatures redraws congressional district boundaries to ensure
that party wins more seats than the voters it represents would otherwise
produce---a practice called gerrymandering.}
In 2019, the U.S.\ Supreme Court held in \textit{Rucho v.\ Common Cause}
that federal courts are powerless to remedy partisan gerrymandering,
leaving Congress as the only institution that can act~\cite{rucho2019}.
Congress has not acted.
The result: in state after state, mapmakers stretch districts into
grotesque shapes, packing opposition voters into a few districts and
cracking communities across many, to predetermine the outcome of
the next decade of elections before a single vote is cast.

\medskip
\noindent\textit{Rucho} explicitly invited Congress to fix the problem.
This brief proposes how.
\noindent\textbf{SUPERSEDED SOURCE --- NOT COMPILED.} The current policy brief
is self-contained in \texttt{main.tex}. This historical section contains
pre-wave claims and must not be reintroduced without a fresh evidence review.
